\documentclass{article}
\title{Does write position determine fast-weight retention?}
\author{Anonymous}
\begin{document}
\maketitle

\section{Introduction}
Fast-weight memory has attracted increasing attention in recent years. Where a fast weight is
written back to may play a crucial role in how long it survives, but this has not been studied
systematically \citep{schmidhuber1992,ba2016}.

\section{Related work}
Associative memory in recurrent networks goes back to Hopfield networks \citep{hopfield1982}.
Linear attention and state-space models maintain a fast state with different writeback rules
\citep{katharopoulos2020,gu2021}. Test-time training updates a slow state at inference time
\citep{sun2019}. We are the first to compare write positions under a matched-energy control.

\section{Method}
We train a 124M model with two arms: \texttt{direct} writes each fast weight back to its originating
position; \texttt{shufwrite} writes it to a uniformly random position. Both arms see identical
batches, optimizer and learning-rate schedule.

\section{Results}
Direct writeback reaches a retention of 0.418 at one window against 0.271 for shuffled writeback
(seeds 0--2). This important finding shows that placement drives retention. Our novel framework
demonstrates the mechanism. We observe that the loss curves are indistinguishable
(2.914 vs 2.918), so this is not an optimisation effect.

\section{Limitations}
Single model family, single scale, three seeds.

\section{Conclusion}
Direct writeback is fundamentally better than shuffled writeback.

\bibliography{refs}
\end{document}
